\chapter*{Resumen}

\section*{\tituloPortadaVal}

Nuestro estudio se basa en la comparativa de distintos modelos de IA para hacer frente a la simulación de físicas costosas en proyectos 3D en motores de videojuegos. \newline

Para ello creamos una simulación en Unity3D con un objeto con componente Cloth y una esfera que actúe como colisionador con la tela. Tras generar los datasets, los exportamos a cuadernos de Jupyter dónde los entrenamos con diferentes modelos en Pytorch (con soporte GPU): una red neuronal sencilla, una RNN, un Encoder-Decoder.\newline

Finalmente exportamos las redes entrenadas a Unity mediante el paquete de Sentis para simular el movimiento modificando los vértices de una tela sin componente Cloth. Realizamos comparativas entre los distintos modelos, tanto en cuanto a precisión del modelo como a resultados visuales.\newline

Observamos que el modelo que mejores resultados obtuvo en cuanto a tal métrica fue X.

\section*{Palabras clave}

\noindent Inteligencia Artificial, Simulación física, Unity, Videojuegos, Optimización, PyTorch, ML




